\label{app:prompts}

This appendix reproduces the prompts, in discharge of guideline~3 of
\textcite{baltes2026llmguidelines}. They are
transcribed from \texttt{prnote/note.py} for generation and from
\texttt{scripts/evaluate\_reviews.py} for scoring, the two files that produced
the reviews and verdicts of Experiments~1 and~2 analysed in
Chapter~\ref{chap:results}. The transcription is character-exact except that
one em dash in the required output heading is set here as a hyphen and the
code fence that wraps the output format in the source is omitted. The human
study used stricter variants of the generator prompts
(Section~\ref{sec:human-study}); these are archived with the study materials
in the replication package rather than reproduced here.

\section{Generator Prompts}
\label{app:prompts-generator}

Every arm shares the following system prompt.

{\footnotesize
\begin{verbatim}
You are a senior software engineer conducting a thorough code review of a
pull request.

Your task is to generate a structured review note that helps the PR author
understand:
1. What potential issues exist
2. Evidence from the code supporting your concerns
3. The impact of these issues
4. Concrete recommendations

Output Format (follow exactly):

# Review Note - Evidence-Anchored

**Scope:** [One sentence describing what this PR does]

## Problem
[Numbered list of 2-3 specific concerns or issues you identified]

## Evidence
[Bullet points with specific file:line references supporting your concerns]

## Impact
[Explain the technical impact - what could go wrong, what risks exist]

## Recommendation (Fix / Tests / Risks)
[Numbered list of concrete, actionable recommendations]

## Traceability
[List relevant code owners or teams if known, otherwise state
"Not specified"]

Guidelines:
- Be specific and cite actual file paths and line numbers from the diff
- Focus on meaningful issues, not style nitpicks
- Consider integration impacts and test coverage
- Be constructive and professional
\end{verbatim}
}

The imposed output structure bears on the measurement. Because all four arms are
told to produce the same five headings, the rubric cannot reward an arm for
organising its output; it can only reward what fills those headings. This narrows
the space in which a difference between arms can appear, and it is one reason the
effects reported in Chapter~\ref{chap:results} are small in absolute terms.

Each arm then appends one paragraph naming the context it has been given. This is
the instruction asymmetry disclosed in Section~\ref{sec:setup}.

\paragraph{\baseline{}.}
{\footnotesize
\begin{verbatim}
Context: You only have access to the PR diff. Make reasonable inferences
about integration and testing needs based on the changes you can see.
\end{verbatim}
}

\paragraph{\kg{}.}
{\footnotesize
\begin{verbatim}
Context: You have access to Knowledge Graph context including:
- Which files are changed and their dependencies
- Which tests cover the changed files
- Who owns the affected code areas
- Which other files call or import the changed code

Use this structural context to provide deeper insights about integration
risks, missing test coverage, and code ownership concerns.
\end{verbatim}
}

\paragraph{\rag{}.}
{\footnotesize
\begin{verbatim}
Context: You have access to semantically similar code chunks from the
repository. These show patterns used elsewhere that may be relevant to
this change.

Use this context to identify consistency issues, suggest similar patterns
to follow, and flag potential conflicts with existing code.
\end{verbatim}
}

\paragraph{\hybrid{}.}
{\footnotesize
\begin{verbatim}
Context: You have access to both:
1. Knowledge Graph context (tests, dependencies, ownership)
2. Semantically similar code from the repository

Combine structural insights (what depends on what, test coverage) with
semantic insights (similar patterns elsewhere) to provide comprehensive
review feedback.
\end{verbatim}
}

The user message is assembled from the frozen evidence pack in a fixed order: the
pull request title; its description, capped at 8\,000 characters and omitted
entirely when empty; the unified diff, capped as given in
Table~\ref{tab:config}; the context block if the arm has one; and the closing
instruction. The diff therefore precedes the context, so in the two arms that
carry a block the added material is the last thing the model reads before the
instruction.
% dataset_v2/scripts/regenerate_reviews_v2.py:110-117 (Experiment 1) and
% prnote/note.py:661-669 (Experiment 2) both order title, body, diff, context.

\section{Per-Criterion Satisfaction Rates}

Table~\ref{tab:exp1-criteria} gives the satisfaction rate behind the
criterion-level analysis. It is here rather than in
Chapter~\ref{chap:results} because no
individual criterion reaches significance once the 17 that move are corrected
for multiplicity, so the chapter's claim is about the band and this table is
reference material for checking it.

\begin{table}[htbp]
    \centering
    \caption{Satisfaction rate per rubric criterion, \baseline{} versus \kg{},
             $n = 35$ (majority of five judges), for all 25 criteria. $\Delta$ is
             in percentage points. The seven criteria marked $\dagger$ sit at $0\,\%$ or
             $100\,\%$ in both arms and cannot register a difference in either
             direction.}
    \label{tab:exp1-criteria}
    \small
    \begin{tabular}{llccc}
        \toprule
        ID & Criterion (abbreviated) & \baseline{} & \kg{} & $\Delta$ \\
        \midrule
        \multicolumn{5}{l}{\emph{KG-relevant band (9 criteria)}} \\
        F3 & Integration with existing code       & 60.0\,\%  & 82.9\,\%  & $+22.9$ \\
        F4 & Breaking changes or impact           & 60.0\,\%  & 74.3\,\%  & $+14.3$ \\
        T1 & Need for tests discussed             & 97.1\,\%  & 100.0\,\% & $+2.9$ \\
        T2 & Edge cases and error paths in tests  & 65.7\,\%  & 68.6\,\%  & $+2.9$ \\
        T3 & Specific test files referenced       & 62.9\,\%  & 71.4\,\%  & $+8.6$ \\
        M1 & Fit with existing architecture       & 8.6\,\%   & 14.3\,\%  & $+5.7$ \\
        M3 & Public API or interface handling     & 8.6\,\%   & 17.1\,\%  & $+8.6$ \\
        C2 & Consistency with similar solutions   & 2.9\,\%   & 5.7\,\%   & $+2.9$ \\
        Q2$^{\dagger}$ & Comments anchored to code locations & 100.0\,\% & 100.0\,\% & $0.0$ \\
        \midrule
        \multicolumn{5}{l}{\emph{Negative-control band (16 criteria)}} \\
        F1 & Change solves the stated problem     & 5.7\,\%   & 8.6\,\%   & $+2.9$ \\
        F2 & Edge cases in the change identified  & 65.7\,\%  & 68.6\,\%  & $+2.9$ \\
        R1 & Code clarity, naming, function size  & 8.6\,\%   & 0.0\,\%   & $-8.6$ \\
        R2 & Unnecessary complexity flagged       & 2.9\,\%   & 2.9\,\%   & $0.0$ \\
        R3 & Comments explain why, not what       & 5.7\,\%   & 2.9\,\%   & $-2.9$ \\
        M2$^{\dagger}$ & Pull request too large to review & 0.0\,\% & 0.0\,\% & $0.0$ \\
        C1 & Project style-guide compliance       & 2.9\,\%   & 0.0\,\%   & $-2.9$ \\
        P1 & Obvious inefficiencies               & 8.6\,\%   & 11.4\,\%  & $+2.9$ \\
        P2$^{\dagger}$ & Benchmarks for critical paths  & 0.0\,\%  & 0.0\,\%  & $0.0$ \\
        S1 & Input validation                     & 14.3\,\%  & 8.6\,\%   & $-5.7$ \\
        S2$^{\dagger}$ & Hardcoded secrets or credentials & 0.0\,\% & 0.0\,\% & $0.0$ \\
        S3 & Error-handling quality               & 20.0\,\%  & 22.9\,\%  & $+2.9$ \\
        Q1$^{\dagger}$ & Overall assessment summary       & 100.0\,\% & 100.0\,\% & $0.0$ \\
        Q3 & Blocking versus non-blocking issues  & 2.9\,\%   & 14.3\,\%  & $+11.4$ \\
        Q4$^{\dagger}$ & Clarifying questions asked       & 0.0\,\% & 0.0\,\% & $0.0$ \\
        Q5$^{\dagger}$ & Rationale given for suggestions  & 100.0\,\% & 100.0\,\% & $0.0$ \\
        \bottomrule
    \end{tabular}
\end{table}

\section{Judge Prompt}
\label{app:prompts-judge}

The panel scores one review against all 25 criteria in a single call. The system
prompt is:

{\footnotesize
\begin{verbatim}
You are an expert code review evaluator. Your task is to assess whether a
PR review meets specific quality criteria.

For each criterion, you must:
1. Determine if the review addresses that criterion (score: 1) or not
   (score: 0)
2. Provide a brief quote or explanation as evidence

Be strict but fair:
- Score 1 only if the review CLEARLY addresses the criterion
- Score 0 if the criterion is not addressed or only vaguely mentioned
- Base your judgment on meaning and intent, not just keywords
\end{verbatim}
}

The user message interpolates the pull request under review, the review text, and
the numbered criteria of Section~\ref{sec:rubric}:

{\footnotesize
\begin{verbatim}
## PR Being Reviewed
Title: {title}
Description: {body, first 400 characters}

## Review to Evaluate
{review, first 4000 characters}

## Criteria to Check
1. [F1] {description}
...
25. [Q5] {description}

## Instructions
Evaluate the review against each criterion. Respond with a JSON object in
this exact format:

{
  "scores": [
    {"id": "F1", "score": 0,
     "evidence": "Review does not verify the change addresses the problem"},
    {"id": "F2", "score": 1,
     "evidence": "Review mentions: 'what about edge cases where X is null'"}
  ]
}

Important:
- Include ALL 25 criteria (F1-F4, T1-T3, R1-R3, M1-M3, C1-C2, P1-P2,
  S1-S3, Q1-Q5)
- Each score must be 0 or 1
- Evidence should be a short quote or explanation
\end{verbatim}
}

Three properties of this prompt bear on how its output should be read.

The two character caps could in principle bias the comparison, since an arm whose
reviews ran longer would have more of its text withheld from the judge. Neither
cap binds in practice. The longest of the 140 canonical reviews is 2\,620
characters and the longest of the 12 human study stimuli is 2\,053, both well
inside the 4\,000-character limit, so no review was truncated before scoring. The
400-character description cap does bind, but it applies identically to all four
arms of a given pull request.

The instruction requires evidence for every verdict, which makes the verdicts
auditable rather than opaque. The stored evidence strings are what
Section~\ref{sec:results-exp1-attribution} inspects when characterising what a
criterion was awarded for.

The judge is asked whether a criterion is \emph{addressed}, not whether the
review is correct about it. This is the coverage-not-correctness limitation that
governs the interpretation of every rubric number in this thesis, stated where
the rubric is defined (Section~\ref{sec:rubric}) and revisited in
Section~\ref{sec:threats}. Reading the prompt shows why the limitation is
structural rather than incidental: a judge given only the review and 400
characters of description has no access to the repository against which
correctness would have to be established.
